\documentclass[12pt,a4paper]{article}
\usepackage[utf8]{inputenc}
\usepackage[T1]{fontenc}
\usepackage[brazil]{babel}
\usepackage[margin=2.5cm]{geometry}
\usepackage{amsmath,amssymb}
\usepackage{booktabs}
\usepackage{graphicx}
\usepackage{float}
\usepackage{tikz}
\usetikzlibrary{arrows.meta,positioning,calc,shapes.geometric}

\newcommand{\Str}{\mathit{Str}}
\newcommand{\Flw}{\mathit{Flw}}
\newcommand{\Li}{\mathit{Li}}
\newcommand{\snow}{\mathit{snow\_covered}}

\title{2\textsuperscript{o} Trabalho de Inteligência Artificial\\[4pt]
\large Raciocínio Probabilístico: rede bayesiana para diagnóstico\\
de um farol de bicicleta movido a dínamo}
\author{Romário Roberto de Pádua Andrade\\
\small Instituto de Computação -- UFAM\\
\small Prof. Edjard Mota}
\date{11 de outubro de 2026}

\begin{document}
\maketitle

\section*{Enunciado resumido}

Variáveis: $\Li$ (luz ligada), $\Str$ (condição da rua: \textit{dry}, \textit{wet},
\textit{snow\_covered}), $\Flw$ (volante do dínamo desgastado), $R$ (dínamo deslizante),
$V$ (dínamo mostra tensão), $B$ (lâmpada ok) e $K$ (cabo ok). Exceto $\Str$, todas
assumem os valores $t/f$.

São dadas as seguintes independências: $\Str$, $\Flw$, $B$ e $K$ são independentes aos
pares; $(R,B)$, $(R,K)$, $(V,B)$ e $(V,K)$ são independentes; e valem
\begin{align*}
P(\Li \mid V, R) &= P(\Li \mid V), \\
P(V \mid R, \Str) &= P(V \mid R), \\
P(V \mid R, \Flw) &= P(V \mid R).
\end{align*}

\section*{(a) Rede bayesiana de causalidade}

% Para usar o seu proprio desenho, comente o ambiente tikzpicture abaixo
% e descomente a linha seguinte (ajustando o nome do arquivo):
% \begin{center}\includegraphics[width=0.75\textwidth]{minha_rede.png}\end{center}

\begin{figure}[H]
\centering
\begin{tikzpicture}[
  no/.style={draw, ellipse, minimum width=1.6cm, minimum height=0.9cm},
  >={Stealth[length=2.5mm]}, node distance=1.4cm and 1.6cm]
  \node[no] (Str) {$\Str$};
  \node[no, right=of Str] (Flw) {$\Flw$};
  \node[no, below=of $(Str)!0.5!(Flw)$] (R) {$R$};
  \node[no, below=of R] (V) {$V$};
  \node[no, right=of V] (B) {$B$};
  \node[no, right=of B] (K) {$K$};
  \node[no, below=of B] (Li) {$\Li$};
  \draw[->] (Str) -- (R);
  \draw[->] (Flw) -- (R);
  \draw[->] (R) -- (V);
  \draw[->] (V) -- (Li);
  \draw[->] (B) -- (Li);
  \draw[->] (K) -- (Li);
\end{tikzpicture}
\caption{Rede bayesiana do farol de bicicleta movido a dínamo.}
\end{figure}

\noindent\textbf{Justificativa das arestas.}
\begin{itemize}
  \item $\Str$, $\Flw$, $B$ e $K$ são independentes aos pares; logo são nós raiz, sem
        arestas entre si.
  \item A condição da rua e o desgaste do volante são as causas do deslizamento do
        dínamo: $\Str \to R$ e $\Flw \to R$.
  \item Como $P(V \mid R, \Str) = P(V \mid R)$ e $P(V \mid R, \Flw) = P(V \mid R)$, o
        único pai de $V$ é $R$: $R \to V$.
  \item A tabela fornecida é $P(\Li \mid V, B, K)$ e, além disso,
        $P(\Li \mid V, R) = P(\Li \mid V)$; logo os pais de $\Li$ são $V$, $B$ e $K$.
  \item As independências $(R,B)$, $(R,K)$, $(V,B)$ e $(V,K)$ garantem que não há
        arestas ligando $B$ ou $K$ a $R$ ou $V$.
\end{itemize}

\section*{(b) e (c) Tabelas de probabilidade condicional}

A única CPT fornecida é a de $\Li$ (Tabela~\ref{tab:li}). Faltam as tabelas de
$\Str$, $\Flw$, $B$, $K$, $R$ e $V$, preenchidas abaixo com valores plausíveis.
Em todas as tabelas, $P(f) = 1 - P(t)$.

\begin{table}[H]
\centering
\begin{tabular}{lc}
\toprule
\multicolumn{2}{c}{Nós raiz} \\
\midrule
$P(\Str = \mathit{dry})$   & 0,70 \\
$P(\Str = \mathit{wet})$   & 0,20 \\
$P(\Str = \snow)$          & 0,10 \\
\midrule
$P(\Flw = t)$              & 0,10 \\
$P(B = t)$                 & 0,95 \\
$P(K = t)$                 & 0,97 \\
\bottomrule
\end{tabular}
\qquad
\begin{tabular}{llc}
\toprule
$\Str$ & $\Flw$ & $P(R = t \mid \Str, \Flw)$ \\
\midrule
\textit{dry}           & $f$ & 0,01 \\
\textit{dry}           & $t$ & 0,20 \\
\textit{wet}           & $f$ & 0,10 \\
\textit{wet}           & $t$ & 0,50 \\
\textit{snow\_covered} & $f$ & 0,40 \\
\textit{snow\_covered} & $t$ & 0,90 \\
\bottomrule
\end{tabular}
\caption{CPTs dos nós raiz e do nó $R$ (valores escolhidos).}
\end{table}

\begin{table}[H]
\centering
\begin{tabular}{cc}
\toprule
$R$ & $P(V = t \mid R)$ \\
\midrule
$t$ & 0,10 \\
$f$ & 0,99 \\
\bottomrule
\end{tabular}
\qquad
\begin{tabular}{cccc}
\toprule
$V$ & $B$ & $K$ & $P(\Li = t \mid V, B, K)$ \\
\midrule
$t$ & $t$ & $t$ & 0,99  \\
$t$ & $t$ & $f$ & 0,01  \\
$t$ & $f$ & $t$ & 0,01  \\
$t$ & $f$ & $f$ & 0,001 \\
$f$ & $t$ & $t$ & 0,3   \\
$f$ & $t$ & $f$ & 0,005 \\
$f$ & $f$ & $t$ & 0,005 \\
$f$ & $f$ & $f$ & 0     \\
\bottomrule
\end{tabular}
\caption{CPT do nó $V$ (valores escolhidos) e CPT do nó $\Li$ (dada no enunciado).}
\label{tab:li}
\end{table}

\noindent\textbf{Critérios para a escolha dos valores.}
\begin{itemize}
  \item Defeitos são eventos raros: $P(\Flw = t)$ é baixo e $P(B = t)$, $P(K = t)$
        são próximos de 1.
  \item A probabilidade de o dínamo deslizar cresce com a piora da rua
        (\textit{dry} $<$ \textit{wet} $<$ \textit{snow\_covered}) e é sempre maior
        com o volante desgastado.
  \item Um dínamo que desliza quase não gera tensão; um dínamo que não desliza gera
        tensão quase sempre.
\end{itemize}

\section*{(d) A rede não contém a aresta $(\Str, \Li)$}

Uma aresta $\Str \to \Li$ só seria necessária se $\Str$ continuasse a influenciar
$\Li$ depois de conhecidos os pais de $\Li$. Isso não ocorre.

\paragraph{Argumento causal.} A condição da rua não age diretamente sobre a luz. Ela
apenas altera a chance de o dínamo deslizar ($R$), o que por sua vez altera a tensão
($V$). Toda a influência passa pelo caminho $\Str \to R \to V \to \Li$.

\paragraph{d-separação.} Esse é o único caminho entre $\Str$ e $\Li$, e ele é uma
cadeia. Em uma cadeia, observar um nó intermediário bloqueia o caminho. Logo, dado
$V$, as variáveis $\Str$ e $\Li$ são condicionalmente independentes:
\[
P(\Li \mid V, \Str) = P(\Li \mid V).
\]

\paragraph{Verificação pelo cálculo.} Como $B$ e $K$ são independentes de $V$ e de
$\Str$,
\[
P(\Li \mid V, \Str) = \sum_{b}\sum_{k} P(\Li \mid V, b, k)\, P(b)\, P(k),
\]
e o lado direito não depende de $\Str$. Para $V = t$:
\begin{align*}
P(\Li = t \mid V = t, \Str)
 &= 0{,}99 \cdot 0{,}95 \cdot 0{,}97 + 0{,}01 \cdot 0{,}95 \cdot 0{,}03 \\
 &\quad + 0{,}01 \cdot 0{,}05 \cdot 0{,}97 + 0{,}001 \cdot 0{,}05 \cdot 0{,}03 \\
 &= 0{,}912285 + 0{,}000285 + 0{,}000485 + 0{,}0000015 \\
 &\approx 0{,}9131 ,
\end{align*}
valor idêntico para $\Str = \mathit{dry}$, $\mathit{wet}$ e $\snow$. Portanto a
aresta $(\Str, \Li)$ não existe na rede.

\paragraph{Observação.} Essa independência condicional é uma hipótese de modelagem,
sustentada pela causalidade do sistema; o cálculo acima mostra que a rede é coerente
com ela.

\section*{(e) Cálculo de $P(V \mid \Str = \snow)$}

Para abreviar, escreve-se $s$ para o evento $\Str = \snow$.
\renewcommand{\snow}{s}

\paragraph{Passo 1: marginalização sobre $R$.}
\[
P(V \mid \snow) = P(V \mid R, \snow)\, P(R \mid \snow)
                + P(V \mid \neg R, \snow)\, P(\neg R \mid \snow).
\]
Usando $P(V \mid R, \Str) = P(V \mid R)$:
\[
P(V \mid \snow) = P(V \mid R)\, P(R \mid \snow) + P(V \mid \neg R)\, P(\neg R \mid \snow).
\]

\paragraph{Passo 2: cálculo de $P(R \mid \snow)$, marginalizando sobre $\Flw$.}
Como $\Flw$ e $\Str$ são independentes, $P(\Flw \mid \snow) = P(\Flw)$. Assim,
\begin{align*}
P(R \mid \snow)
 &= P(R \mid \snow, \Flw)\, P(\Flw) + P(R \mid \snow, \neg \Flw)\, P(\neg \Flw) \\
 &= 0{,}90 \cdot 0{,}10 + 0{,}40 \cdot 0{,}90 \\
 &= 0{,}09 + 0{,}36 = 0{,}45 ,
\end{align*}
e, portanto, $P(\neg R \mid \snow) = 1 - 0{,}45 = 0{,}55$.

\paragraph{Passo 3: substituição.}
\begin{align*}
P(V \mid \snow) &= 0{,}10 \cdot 0{,}45 + 0{,}99 \cdot 0{,}55 \\
                &= 0{,}045 + 0{,}5445 = 0{,}5895 .
\end{align*}

\paragraph{Resposta.} $P(V = t \mid \Str = \mathit{snow\_covered}) \approx 0{,}59$. Com a rua coberta de
neve, o dínamo gera tensão em cerca de 59\% das vezes.

\end{document}
